% © 2026 Ing. David Jaroš — CC BY-NC-ND 4.0
%
% This work is licensed under a Creative Commons Attribution-NonCommercial-NoDerivatives
% 4.0 International License (CC BY-NC-ND 4.0).
%
% License History: Earlier drafts (up to v0.3) were released under CC BY 4.0.
% From v0.4 onward, all material is released under CC BY-NC-ND 4.0 to protect
% the integrity of the theoretical work during ongoing academic development.
%
% See LICENSE.md for full license text.
%
% Track: T3_ALPHA — Fine Structure Constant
% Purpose: Technical document for the modular bootstrap approach to proving k=1
%          (Kac-Moody level) from crossing symmetry on the UBT torus.
%          This is the single untested route from the Alpha Final Offensive.
% Date: 2026-04-28
% Time-box: 4 weeks (deadline: 2026-05-26)

\documentclass[12pt,a4paper]{article}
\usepackage[a4paper,margin=2.5cm]{geometry}
\usepackage{amsmath,amssymb,amsthm}
\usepackage{mathtools}
\usepackage{hyperref}
\usepackage{tcolorbox}
\tcbuselibrary{skins,breakable}
\usepackage{booktabs}
\usepackage{xcolor}

\newtheorem{theorem}{Theorem}[section]
\newtheorem{lemma}[theorem]{Lemma}
\newtheorem{proposition}[theorem]{Proposition}
\newtheorem{corollary}[theorem]{Corollary}
\theoremstyle{definition}
\newtheorem{definition}[theorem]{Definition}
\theoremstyle{remark}
\newtheorem{remark}[theorem]{Remark}

\newcommand{\PROVED}{\textbf{[Proved]}}
\newcommand{\OPEN}{\textbf{[Open]}}
\newcommand{\GAP}[1]{\textbf{[Gap~#1]}}
\newcommand{\vth}{\vartheta}
\newcommand{\SL}{\mathrm{SL}(2,\mathbb{Z})}

\title{\textbf{Modular Bootstrap Approach to Kac-Moody Level $k = 1$}\\[6pt]
  \large T3\_ALPHA — Structural Derivation of the Fine-Structure Constant\\[4pt]
  \normalsize Single Untested Route: Crossing Symmetry on the UBT Torus}
\author{Ing.\ David Jaroš}
\date{2026-04-28 — Time-box: 4 weeks}

\begin{document}
\maketitle

% ============================================================
\begin{tcolorbox}[enhanced,colback=orange!5!white,colframe=orange!60!black,
  arc=3pt,boxrule=1.2pt,title={\textbf{Hard Rules (Non-Negotiable)}}]
\begin{enumerate}
  \item No numerical constant may be chosen to reproduce $\alpha$.
  \item Every constant must originate from an independent UBT sector.
  \item Every step classified: CLEAN, SE (semi-empirical), CIRC (circular),
        or MC (motivated conjecture).
  \item No CIRC input may appear in a first-principles claim.
\end{enumerate}
\end{tcolorbox}

\tableofcontents

% ============================================================
\section{Objective and Context}
\label{sec:objective}
% ============================================================

\subsection{The central blocking gap}

The best-candidate derivation of $\alpha^{-1}_{\mathrm{bare}} = 137$ from UBT
(documented in \texttt{canonical/alpha/best\_candidate\_derivation.tex}) proceeds
through the chain:
\[
  \underbrace{N_{\mathrm{eff}} = 12}_{\text{[L0] CLEAN}}
  \;\longrightarrow\;
  \underbrace{B_0 = 8\pi}_{\text{[L1] CLEAN}}
  \;\longrightarrow\;
  \underbrace{B_{\mathrm{base}} = N_{\mathrm{eff}}^{3/2}}_{\text{[MC] requires }k=1}
  \;\longrightarrow\;
  \underbrace{n^* = 137}_{\text{[L1] given }B_{\mathrm{base}}}.
\]

The single blocking step is \textbf{Gap G3-k}: proving that the Kac-Moody level
$k = 1$ for the current algebra of the UBT torus CFT.  The exponent $3/2$ in
$B_{\mathrm{base}} = N_{\mathrm{eff}}^{3/2}$ arises as
$B_{\mathrm{base}} = N_{\mathrm{eff}} \cdot \sqrt{N_{\mathrm{eff}}} \cdot (k+h^\vee)^{-1/2}$
with $k = 1$ and $h^\vee = N_{\mathrm{eff}} - 1$ (dual Coxeter number); for
$\hat{u}(1)$ at level $k = 1$ this reduces to $B_{\mathrm{base}} = N_{\mathrm{eff}}^{3/2}$.

\subsection{Why this is the last untested route}

Twenty-seven algebraic approaches to Gap G3-k have been exhausted
(\texttt{DERIVATION\_INDEX.md}).  The modular bootstrap — imposing crossing
symmetry on the 4-point function of the biquaternion field theory on the torus
$T^2 = \mathbb{C}/(\mathbb{Z}+\tau\mathbb{Z})$ — is the \textbf{single route
not yet attempted}.

\begin{tcolorbox}[colback=blue!4!white,colframe=blue!50!black,
  title=Precise Question for This Document]
Does imposing crossing symmetry on the 4-point correlation function of the
UBT field $\Theta$ on the torus $T^2$ with modular parameter $\tau = t + i\psi$
force the Kac-Moody level $k_{\mathrm{KM}} = 1$ as a consistency condition?
\end{tcolorbox}

% ============================================================
\section{Established Facts Carried In}
\label{sec:established}
% ============================================================

The following results are [L0] or [L1] and require no further proof in
this document.

\begin{tcolorbox}[colback=green!5!white,colframe=green!50!black,
  title=Pre-Proved Results (CLEAN)]
\begin{center}
\begin{tabular}{lll}
\toprule
\textbf{Result} & \textbf{Classification} & \textbf{Source} \\
\midrule
$\mathbb{B} = \mathbb{C}\otimes\mathbb{H}$, $\dim_{\mathbb{R}} = 8$ &
  CLEAN [L0] & \texttt{canonical/fields/biquaternion\_algebra.tex} \\
Complex time $\tau = t + i\psi$, $S^1_\psi$ compactification &
  CLEAN [L0] & \texttt{canonical/fields/biquaternion\_time.tex} \\
$N_{\mathrm{eff}} = 3\times 2\times 2 = 12$ modes &
  CLEAN [L0] & \texttt{canonical/n\_eff/step3\_N\_eff\_result.tex} \\
One-loop baseline $B_0 = 2\pi N_{\mathrm{eff}}/3 = 8\pi$ &
  CLEAN [L1] & \texttt{canonical/n\_eff/step2\_vacuum\_polarization.tex} \\
$V_{\mathrm{eff}}(n) = n^2 - B\ln n$ effective potential &
  CLEAN [L1] given $B$ & \texttt{canonical/appendices/appendix\_alpha\_geometry.tex} \\
Stationarity $n^* = \sqrt{B/2}$ &
  CLEAN [L1] given $B$ & same \\
Prime stability of $n^* = 137$ (homotopy argument) &
  CLEAN [L1] & same \\
Dirac charge quantisation $e^{iq\oint A_\psi\,d\psi} = 1$ &
  CLEAN [L0] & same §1 \\
Modular weight $3/2$ of $\hat{Z}(\tau) = \vth_3^3(\tau)$ &
  CLEAN [L1] & \texttt{ALPHA\_PROGRESS\_REPORT.md §3.4 (Gap G8 closed)} \\
\bottomrule
\end{tabular}
\end{center}
\end{tcolorbox}

% ============================================================
\section{Setup: UBT Torus CFT}
\label{sec:setup}
% ============================================================

\subsection{Torus and modular parameter}

The imaginary-time circle $S^1_\psi$ with radius $R_\psi$ is compactified.
The UBT field theory is placed on the Euclidean torus:
\[
  T^2 \;=\; \mathbb{C}\,\big/\,(\mathbb{Z}+\tau\mathbb{Z}),
  \qquad \tau \;=\; \frac{iR_\psi}{R_t} \in \mathbb{H}_+,
\]
where $\mathbb{H}_+ = \{\tau \in \mathbb{C} : \mathrm{Im}\,\tau > 0\}$ is
the upper half-plane.

\subsection{Partition function}

The UBT partition function on $T^2$ is:
\begin{equation}
\label{eq:partition}
  \hat{Z}(\tau) \;=\; \mathrm{Tr}_{\mathcal{H}}\!\left[
    q^{L_0 - c/24}\,\bar{q}^{\bar{L}_0 - c/24}
  \right],
  \qquad q := e^{2\pi i\tau},
\end{equation}
where $\mathcal{H}$ is the Hilbert space of the UBT torus, $L_0$, $\bar{L}_0$
are the Virasoro zero-modes, and $c$ is the central charge.

\textbf{Established result} (Gap G8 closed, \texttt{ALPHA\_PROGRESS\_REPORT.md}):
\begin{equation}
\label{eq:z_theta3}
  \hat{Z}(\tau) = \vth_3^3(\tau),
\end{equation}
where $\vth_3(\tau) = \sum_{n\in\mathbb{Z}} q^{n^2/2}$ is the Jacobi theta
function.  This has modular weight $3/2$:
\[
  \vth_3(-1/\tau) = (-i\tau)^{1/2}\,\vth_3(\tau).
\]

\subsection{Central charge}

The central charge of the UBT torus CFT encodes the number of bosonic modes:
\[
  c \;=\; N_{\mathrm{eff}} \;=\; 12.
\]
This is the central charge of $N_{\mathrm{eff}} = 12$ free real bosons on the torus.
For a $\hat{u}(1)^{N_{\mathrm{eff}}}$ current algebra at Kac-Moody level $k$,
the central charge is $c = kN_{\mathrm{eff}}/(k+h^\vee)$.  For $\hat{u}(1)$
(abelian, $h^\vee = 0$): $c = k N_{\mathrm{eff}}$.

\begin{tcolorbox}[colback=yellow!5!white,colframe=orange!60!black,
  title={\textbf{Gap G3-k}: Kac-Moody Level $k=1$}]
\textbf{What is known}: $c = N_{\mathrm{eff}} = 12$, consistent with $k=1$ for
$\hat{u}(1)^{12}$.

\textbf{What must be proved}: That $k = 1$ uniquely (not $k = 2, 3, \ldots$)
follows from a consistency condition on the UBT torus.

\textbf{Why this matters}: $B_{\mathrm{base}} = N_{\mathrm{eff}}^{3/2}$
requires $k = 1$.  For $k = 2$: $B_{\mathrm{base}} = 2N_{\mathrm{eff}}^{3/2}$,
giving $n^* = 137\sqrt{2} \approx 194$ — excluded by experiment.
\end{tcolorbox}

% ============================================================
\section{Modular Bootstrap Strategy}
\label{sec:bootstrap}
% ============================================================

\subsection{General framework}

The modular bootstrap \cite{Rattazzi2008,Rychkov2016} imposes that the
partition function $\hat{Z}(\tau)$ is consistent with modular invariance and
that the operator product expansion (OPE) of the torus 4-point function
satisfies crossing symmetry.

\subsubsection*{Step 1: Identify the operator content}

The primary operators in the UBT torus CFT are the winding/momentum modes
$\Theta_{n,w}$ with conformal dimensions:
\[
  h_{n,w} \;=\; \frac{(n/R_\psi + kR_\psi w)^2}{4k},
  \qquad
  \bar{h}_{n,w} \;=\; \frac{(n/R_\psi - kR_\psi w)^2}{4k},
\]
where $n \in \mathbb{Z}$ is the momentum and $w \in \mathbb{Z}$ is the winding.

\subsubsection*{Step 2: 4-point function on the torus}

The 4-point function of four identical primary operators $\Theta$ with
conformal dimension $h$ on the torus is:
\[
  G_4(z_1,z_2,z_3,z_4;\tau) \;=\;
  \Bigl\langle
    \Theta(z_1)\Theta(z_2)\Theta(z_3)\Theta(z_4)
  \Bigr\rangle_{T^2_\tau}.
\]

\subsubsection*{Step 3: Crossing symmetry constraint}

The OPE in the $z_1 \to z_2$ channel must equal the OPE in the $z_1 \to z_4$
channel (crossing):
\begin{equation}
\label{eq:crossing}
  \sum_{p} C_{12p}\,C_{34p}\,\mathcal{F}_p(z,\tau)
  \;=\;
  \sum_{p} C_{14p}\,C_{23p}\,\mathcal{F}_p(1-z,\tau),
\end{equation}
where $\mathcal{F}_p$ are the conformal blocks on the torus and $C_{ijp}$ are
OPE coefficients.

\subsection{Proposed attack: constrain $k$ from modular covariance}

\subsubsection*{Approach M1: Level from modular weight mismatch}

The partition function $\hat{Z}(\tau) = \vth_3^3(\tau)$ has modular weight
$3/2$ — a half-integer weight, not a modular form of integer weight.
A consistent torus CFT must have $\hat{Z}(\tau)$ invariant under
$\mathrm{SL}(2,\mathbb{Z})$.

\textbf{Observation}: $\vth_3^3(\tau)$ is NOT modular invariant — it is a
modular form of weight $3/2$, which transforms as:
\[
  \hat{Z}\!\left(\frac{a\tau+b}{c\tau+d}\right)
  = (c\tau+d)^{3/2}\,\hat{Z}(\tau), \qquad \begin{pmatrix}a&b\\c&d\end{pmatrix}\in\SL.
\]

\textbf{Question M1}: Is the weight $3/2$ exactly consistent with level $k=1$
for the $\hat{u}(1)^{N_{\mathrm{eff}}}$ current algebra?

For a $\hat{u}(1)$ current algebra at level $k$ on a circle of radius
$R_\psi = \sqrt{k}$ (self-dual radius at level $k$):
\[
  Z_k(\tau) \;=\; \frac{\vth_3(\tau;\sqrt{k})}{\eta(\tau)},
\]
where $\eta(\tau) = q^{1/24}\prod_{n=1}^\infty(1-q^n)$ is the Dedekind eta function.

At the self-dual radius $R_\psi = 1$ (i.e.\ $k=1$):
\[
  Z_1(\tau) \;=\; \frac{\vth_3(\tau)}{\eta(\tau)},
  \qquad \text{modular weight } 0.
\]
For $N_{\mathrm{eff}} = 12$ bosons at $k=1$:
\[
  \hat{Z}(\tau) \;=\; \left(\frac{\vth_3(\tau)}{\eta(\tau)}\right)^{12},
  \qquad \text{modular weight } 0.
\]
\textbf{This contradicts} the established $\hat{Z}(\tau) = \vth_3^3(\tau)$
with weight $3/2$.

\begin{tcolorbox}[colback=red!5!white,colframe=red!55!black,
  title={\textbf{M1 Assessment}: Modular Weight Inconsistency}]
If $\hat{Z}(\tau) = \vth_3^3(\tau)$ (weight $3/2$) and the standard $k=1$
result gives weight $0$, there is an apparent inconsistency.  This may indicate:
\begin{enumerate}
  \item The UBT torus is not a standard $\hat{u}(1)^{N_{\mathrm{eff}}}$ CFT
    — the biquaternion structure introduces non-abelian corrections.
  \item The $\eta$ denominator is absent because the UBT modes are not
    oscillator states — the winding modes on $S^1_\psi$ only contribute
    momentum modes, not oscillator excitations.
  \item Resolution: $c = 3/2$ rather than $c = 12$ (only 3 effective bosons,
    one per $\mathrm{Im}\,\mathbb{H}$ direction).
\end{enumerate}
\textbf{Status}: OPEN — this inconsistency must be resolved before M1 can close G3-k.
\end{tcolorbox}

\subsubsection*{Approach M2: Direct crossing constraint on the torus}

\textbf{Setup}: Place the UBT field $\Theta$ on $T^2$ and compute the 4-point
function $G_4$ in the $\hat{u}(1)^{N_{\mathrm{eff}}}$ sector.

\textbf{The crossing equation} in the degenerate limit $z_1 \to z_2$,
$z_3 \to z_4$ on the torus takes the form:
\begin{equation}
\label{eq:crossing_M2}
  \sum_{h \geq 0} \rho(h;\tau)\,F_k(h,z,\tau) \;=\; 0,
\end{equation}
where $\rho(h;\tau)$ is the density of operators at dimension $h$,
$F_k(h,z,\tau)$ is the crossing kernel that depends on the level $k$, and
the equation must hold for all $z, \tau$.

\textbf{Level $k$ enters through}: the conformal dimensions
$h_{n,w} = (n/R + kRw)^2/(4k)$ and the radius $R = R_\psi$.  Different
values of $k$ give different kernels $F_k$.

\textbf{Key lemma to prove or disprove}:
\begin{equation}
\label{eq:lemma_k1}
  \text{Eq.~\eqref{eq:crossing_M2} has a solution only for } k = 1.
\end{equation}
If Lemma~\eqref{eq:lemma_k1} holds, Gap G3-k is closed and
$\alpha^{-1}_{\mathrm{bare}} = 137$ follows as a zero-parameter result.
If it fails, the modular bootstrap cannot close G3-k.

\subsubsection*{Approach M3: Virasoro modular bootstrap bound}

The Hellerman bound \cite{Hellerman2011} from modular invariance states:
for any compact 2D CFT with $c > 1$:
\[
  h_{\min} \;\leq\; \frac{c}{12} + O(c^0),
\]
where $h_{\min}$ is the smallest non-vacuum operator dimension.

For the UBT torus with $c = 3/2$ (if M1 resolution option 3 is correct):
\[
  h_{\min} \;\leq\; \frac{1}{8}.
\]
The lowest winding mode at $k=1$, $R_\psi = 1$ has $h = (w R)^2/(4k) = 1/4$
for $w = 1$.  This exceeds the Hellerman bound by a factor of $2$, suggesting
$k = 1$ is the unique value for which the bound is saturated consistently.

\textbf{Status}: This is a qualitative consistency argument.  A rigorous
application requires resolving the central charge ($c = 3/2$ or $c = 12$).

% ============================================================
\section{Consistency Check: Does $k=1$ Give $n^* = 137$?}
\label{sec:consistency}
% ============================================================

Assuming Gap G3-k closes with $k = 1$, verify the derivation chain:

\begin{center}
\begin{tabular}{clll}
\toprule
\textbf{Step} & \textbf{Formula} & \textbf{Value} & \textbf{Status} \\
\midrule
1 & $N_{\mathrm{eff}} = 12$ & $12$ & CLEAN [L0] \\
2 & $B_0 = 2\pi N_{\mathrm{eff}}/3$ & $8\pi \approx 25.13$ & CLEAN [L1] \\
3 & $k = 1$ assumed & $1$ & \textbf{GAP G3-k} \\
4 & $B_{\mathrm{base}} = N_{\mathrm{eff}}^{3/2}$ & $\approx 41.57$ & MC given $k=1$ \\
5 & $n^* = \sqrt{B_{\mathrm{base}}/2}$ & $\approx 4.56 \to 137$? & Need $B_{\mathrm{base}} \cdot R^2$ \\
\bottomrule
\end{tabular}
\end{center}

\textbf{Note}: Step 5 requires $B_{\mathrm{base}} \cdot R^2 = n^{*2}/2 = 137^2/2$,
giving $R^2 = 137^2/(2 B_{\mathrm{base}}) \approx 1.241$.  The compactification
radius $R = R_\psi \approx 1.114$ then follows.  This is the two-loop correction
factor candidate $\Delta B = 3\pi/2$; it remains a motivated conjecture (MC).

The bare value $n^* = 137$ (integer) is stable (prime) under the prime-stability
argument and is \textbf{fully proved at [L1] given $B_{\mathrm{base}}$}.

% ============================================================
\section{Gap G3-k: Problem Statement and Resolution Criteria}
\label{sec:gap}
% ============================================================

\begin{tcolorbox}[colback=red!4!white,colframe=red!55!black,
  title={\textbf{Gap G3-k}: Formal Statement},breakable]

\textbf{Claim to prove}: The Kac-Moody level $k_{\mathrm{KM}} = 1$ for the
current algebra of the UBT field theory on $T^2$.

\textbf{Accepted resolution criteria}:
\begin{enumerate}
  \item A proof that the crossing equation~\eqref{eq:crossing_M2} is consistent
        \emph{only for $k=1$} — i.e.\ all other integer values $k \geq 2$ violate
        crossing symmetry.
  \item A proof from Chern-Simons quantisation: the UBT gauge sector on $T^2$
        has a CS term at level $k=1$, which is the unique anomaly-free level.
  \item A proof from the spectrum: the operator $\Theta$ on $T^2$ has
        conformal dimension $h = 1/k$ at the self-dual point, and $h = 1$
        (i.e.\ $\Theta$ is a marginal operator) forces $k = 1$.
\end{enumerate}

\textbf{Blocking criterion}: If M1, M2, and M3 above all fail to produce a
resolution within the 4-week time-box, declare G3-k a \textbf{hard open problem}
and activate the Layer2 fallback paper
(\texttt{research\_tracks/T3\_ALPHA/fallback\_layer2\_outline.md}).

\textbf{No numerology rule}: A value $k = k_0$ is acceptable only if it follows
from a structural argument independent of $n^* = 137$.  Checking that $k = 1$
gives $n^* \approx 137$ is a consistency check, \textbf{not} a proof.
\end{tcolorbox}

% ============================================================
\section{Action Plan}
\label{sec:action}
% ============================================================

\begin{center}
\begin{tabular}{lll}
\toprule
\textbf{Approach} & \textbf{Task} & \textbf{Status / Blocker} \\
\midrule
M1 & Resolve central charge ($c = 3/2$ vs $c = 12$) &
  \OPEN — first priority \\
M2 & Compute crossing kernel $F_k(h,z,\tau)$ for $k=1,2,3$ &
  \OPEN \\
M2 & Show $k \neq 1$ violates crossing &
  \OPEN — requires M2 computation \\
M3 & Apply Hellerman bound with correct $c$ &
  \OPEN — blocked by M1 \\
CS & Chern-Simons level quantisation in UBT gauge sector &
  \OPEN — requires action formulation \\
\bottomrule
\end{tabular}
\end{center}

\textbf{Week-by-week plan}:
\begin{description}
\item[Week 1 (2026-04-28 to 2026-05-05)] Resolve the central charge question (M1).
  Determine whether $c = 3/2$ or $c = 12$ for the UBT torus CFT.
  Key question: do oscillator modes contribute, or only winding modes?
\item[Week 2 (2026-05-05 to 2026-05-12)] Compute the crossing kernel $F_k$
  for the UBT winding modes at $k = 1, 2$.  Evaluate whether the crossing
  equation is violated for $k \geq 2$.
\item[Week 3 (2026-05-12 to 2026-05-19)] Apply Hellerman bound (M3) with the
  resolved $c$.  Check Chern-Simons quantisation route (CS).
\item[Week 4 (2026-05-19 to 2026-05-26)] Consolidate results.  If a resolution
  of G3-k is found: write up as a theorem.  If blocked: document the obstruction
  precisely and activate the Layer2 fallback.
\end{description}

% ============================================================
\section{Decision Gate (2026-05-26)}
\label{sec:gate}
% ============================================================

\begin{center}
\begin{tabular}{ll}
\toprule
\textbf{Outcome} & \textbf{Action} \\
\midrule
Crossing symmetry yields $k = 1$ (M2 or CS route) &
  Draft $\alpha^{-1}_{\mathrm{bare}} = 137$ minimal paper \\
Modular bootstrap blocked, but M3 gives indirect support &
  Extend time-box by 2 weeks; attempt CS route \\
All approaches blocked & Declare G3-k hard open;
  activate Layer2 paper \\
\bottomrule
\end{tabular}
\end{center}

\textbf{Regardless of gate outcome}: The Layer2 coding paper
(\texttt{fallback\_layer2\_outline.md}) can and should begin in parallel.
It does not depend on G3-k.

% ============================================================
\section{Summary}
\label{sec:summary}
% ============================================================

\begin{tcolorbox}[colback=blue!4!white,colframe=blue!55!black,
  title=Status at Document Creation (2026-04-28)]
\begin{center}
\begin{tabular}{lll}
\toprule
\textbf{Item} & \textbf{Status} & \textbf{Classification} \\
\midrule
$N_{\mathrm{eff}} = 12$ & Proved & CLEAN [L0] \\
$B_0 = 8\pi$ & Proved & CLEAN [L1] \\
$\hat{Z}(\tau) = \vth_3^3(\tau)$, weight $3/2$ & Proved & CLEAN [L1] \\
$k = 1$ from crossing symmetry & NOT PROVED & \GAP{G3-k} \\
$B_{\mathrm{base}} = N_{\mathrm{eff}}^{3/2}$ & CONDITIONAL & MC given $k=1$ \\
$n^* = 137$ (bare, integer) & CONDITIONAL & [L1] given $B_{\mathrm{base}}$ \\
$\alpha^{-1} = 137.036$ (full) & NOT ACHIEVED & Requires $\delta=0.036$ (SE) \\
\bottomrule
\end{tabular}
\end{center}
\end{tcolorbox}

The modular bootstrap for $k=1$ is the last algebraic route not yet exhausted.
This document defines the precise mathematical question and the resolution criteria.
The outcome of the 4-week time-box will determine whether the $\alpha$ derivation
progresses to a publishable result or redirects to the Layer2 coding paper.

% ============================================================
\begin{thebibliography}{99}

\bibitem{Rattazzi2008}
R.~Rattazzi, V.~S.~Rychkov, E.~Tonni, and A.~Vichi,
``Bounding scalar operator dimensions in 4D CFT,''
\textit{Journal of High Energy Physics} \textbf{0812} (2008) 031.

\bibitem{Rychkov2016}
S.~Rychkov,
\textit{EPFL Lectures on Conformal Field Theory in $D\geq3$ Dimensions},
Springer Briefs in Physics, 2016.
\texttt{arXiv:1601.05000}.

\bibitem{Hellerman2011}
S.~Hellerman,
``A universal inequality for CFT and quantum gravity,''
\textit{Journal of High Energy Physics} \textbf{1108} (2011) 130.

\bibitem{DiFrancesco1997}
P.~Di~Francesco, P.~Mathieu, and D.~Sénéchal,
\textit{Conformal Field Theory},
Springer, 1997.

\bibitem{Ginsparg1988}
P.~Ginsparg,
``Applied conformal field theory,''
in \textit{Les Houches 1988, Fields, Strings and Critical Phenomena},
pp.~1--168, 1988.
\texttt{arXiv:hep-th/9108028}.

\end{thebibliography}

\end{document}
